\chapter*{Conclusion}
\addcontentsline{toc}{chapter}{Conclusion}

In the end we managed to develop two nontrivial applications. One for running and simulating the user's Arduino code. The other one for creating a interactive virtual environment, where the Arduino can be connected into a circuit with visual feedback. In chapter \ref{ch:architecture} of this thesis we described the thought process, project structure, and technologies used in order to achieve the goals we set for ourselves in section \ref{sec:goals}. Then we looked into more detail on how we implemented these technologies while avoiding many pitfalls in chapters \ref{ch:backend} and \ref{ch:frontend}.

While we implemented only a couple of most used functions from the Arduino library, and the user interface leaves a lot to be desired, we believe the current version of the product represents a solid foundation on which new functions, events, components and quality of life features can be laid down. One of the biggest hurdles was making sure that the two applications stay synchronized. We succeeded in that regard, by learning a lot of useful information about interprocess communication and approaches commonly used in distributed computing.

A big success also lies in the extensive component system developed in the frontend. The user can extend their component library just by following the specification in section \ref{sec:component_spec}. They can use custom vector graphics, and give life to their components thanks to the abstract \texttt{Component} class. Then they can add it into a scene together with Arduino running their own code. All this while all the heavy work is done by the reflection and parsing machinery on the background. We would be excited to work on this component system, so that it brings even more possibilities to the user.